\documentclass[11pt]{article}
\newcommand{\doctitle}{Two Architectures for Simulating
  Biomarker-Treatment Interactions: Implications for Statistical
  Power Under Carryover}
\newcommand{\shorttitle}{DGP Architecture Comparison}
\newcommand{\docdate}{2026-04-09}
% pmsimstats-ng standard document preamble
% Usage: % pmsimstats-ng standard document preamble
% Usage: % pmsimstats-ng standard document preamble
% Usage: \input{pmsimstats-preamble} after \documentclass[11pt]{article}
%
% Documents must define before \input:
%   \newcommand{\doctitle}{...}       Full document title
%   \newcommand{\shorttitle}{...}     Short title for header
%   \newcommand{\docdate}{YYYY-MM-DD} Document date

\usepackage[margin=1in]{geometry}
\usepackage{amsmath,amssymb}
\usepackage[T1]{fontenc}
\usepackage{lmodern}
\usepackage{microtype}
\usepackage{graphicx}
\graphicspath{{figures/}}
\usepackage{booktabs}
\usepackage{longtable}
\usepackage{array}
\usepackage{hyperref}
\usepackage{xcolor}
\usepackage{fancyhdr}
\usepackage{parskip}
\usepackage{caption}
\usepackage{enumitem}
\usepackage{verbatim}
\usepackage{listings}

\lstset{
  language=R,
  basicstyle=\small\ttfamily,
  keywordstyle=\color{blue!70!black},
  commentstyle=\color{green!50!black},
  stringstyle=\color{red!60!black},
  breaklines=true,
  frame=single,
  framerule=0.5pt,
  rulecolor=\color{gray!40},
  backgroundcolor=\color{gray!5},
  xleftmargin=1em,
  framexleftmargin=1em,
  columns=flexible
}

\hypersetup{
  colorlinks=true,
  linkcolor=blue!60!black,
  citecolor=green!50!black,
  urlcolor=blue!60!black
}

\pagestyle{fancy}
\fancyhf{}
\lhead{\footnotesize pmsimstats-ng Technical Report}
\rhead{\footnotesize \shorttitle}
\rfoot{\footnotesize \thepage}
\renewcommand{\headrulewidth}{0.4pt}

\title{\doctitle}
\author{pmsimstats team}
\date{\docdate}
 after \documentclass[11pt]{article}
%
% Documents must define before \input:
%   \newcommand{\doctitle}{...}       Full document title
%   \newcommand{\shorttitle}{...}     Short title for header
%   \newcommand{\docdate}{YYYY-MM-DD} Document date

\usepackage[margin=1in]{geometry}
\usepackage{amsmath,amssymb}
\usepackage[T1]{fontenc}
\usepackage{lmodern}
\usepackage{microtype}
\usepackage{graphicx}
\graphicspath{{figures/}}
\usepackage{booktabs}
\usepackage{longtable}
\usepackage{array}
\usepackage{hyperref}
\usepackage{xcolor}
\usepackage{fancyhdr}
\usepackage{parskip}
\usepackage{caption}
\usepackage{enumitem}
\usepackage{verbatim}
\usepackage{listings}

\lstset{
  language=R,
  basicstyle=\small\ttfamily,
  keywordstyle=\color{blue!70!black},
  commentstyle=\color{green!50!black},
  stringstyle=\color{red!60!black},
  breaklines=true,
  frame=single,
  framerule=0.5pt,
  rulecolor=\color{gray!40},
  backgroundcolor=\color{gray!5},
  xleftmargin=1em,
  framexleftmargin=1em,
  columns=flexible
}

\hypersetup{
  colorlinks=true,
  linkcolor=blue!60!black,
  citecolor=green!50!black,
  urlcolor=blue!60!black
}

\pagestyle{fancy}
\fancyhf{}
\lhead{\footnotesize pmsimstats-ng Technical Report}
\rhead{\footnotesize \shorttitle}
\rfoot{\footnotesize \thepage}
\renewcommand{\headrulewidth}{0.4pt}

\title{\doctitle}
\author{pmsimstats team}
\date{\docdate}
 after \documentclass[11pt]{article}
%
% Documents must define before \input:
%   \newcommand{\doctitle}{...}       Full document title
%   \newcommand{\shorttitle}{...}     Short title for header
%   \newcommand{\docdate}{YYYY-MM-DD} Document date

\usepackage[margin=1in]{geometry}
\usepackage{amsmath,amssymb}
\usepackage[T1]{fontenc}
\usepackage{lmodern}
\usepackage{microtype}
\usepackage{graphicx}
\graphicspath{{figures/}}
\usepackage{booktabs}
\usepackage{longtable}
\usepackage{array}
\usepackage{hyperref}
\usepackage{xcolor}
\usepackage{fancyhdr}
\usepackage{parskip}
\usepackage{caption}
\usepackage{enumitem}
\usepackage{verbatim}
\usepackage{listings}

\lstset{
  language=R,
  basicstyle=\small\ttfamily,
  keywordstyle=\color{blue!70!black},
  commentstyle=\color{green!50!black},
  stringstyle=\color{red!60!black},
  breaklines=true,
  frame=single,
  framerule=0.5pt,
  rulecolor=\color{gray!40},
  backgroundcolor=\color{gray!5},
  xleftmargin=1em,
  framexleftmargin=1em,
  columns=flexible
}

\hypersetup{
  colorlinks=true,
  linkcolor=blue!60!black,
  citecolor=green!50!black,
  urlcolor=blue!60!black
}

\pagestyle{fancy}
\fancyhf{}
\lhead{\footnotesize pmsimstats-ng Technical Report}
\rhead{\footnotesize \shorttitle}
\rfoot{\footnotesize \thepage}
\renewcommand{\headrulewidth}{0.4pt}

\title{\doctitle}
\author{pmsimstats team}
\date{\docdate}


\begin{document}
\maketitle
\thispagestyle{fancy}
\tableofcontents

\bigskip

\textbf{pmsimstats team}

\bigskip\noindent\rule{\textwidth}{0.4pt}\bigskip

\section{Abstract}

Simulation studies of predictive biomarker-treatment interactions
in clinical trials require a data-generating process (DGP) that
specifies how the biomarker moderates treatment response. We
identify two fundamentally different DGP architectures used in
practice and show that the choice between them determines the
extent to which carryover effects reduce statistical power to
detect the interaction. Under \textit{direct mean moderation}, the
biomarker scales the treatment effect in the population mean
structure; power is largely preserved under carryover because the
proportional relationship between biomarker and drug response is
maintained at every timepoint. Under \textit{differential
correlation} (the multivariate normal, or MVN, approach), the interaction emerges from treatment-state-dependent
correlation between biomarker and response in the joint
distribution; carryover erodes this differential correlation signal
and substantially reduces power. We formalize the mathematical
distinction, illustrate the consequences with simulation results
from an N-of-1 trial design framework, discuss the biological
assumptions underlying each architecture, and provide guidance for
selecting the appropriate DGP for a given clinical context.

\bigskip\noindent\rule{\textwidth}{0.4pt}\bigskip

\section{1. Introduction}

A central question in precision medicine trial design is whether a
baseline biomarker predicts differential treatment response -- that
is, whether the treatment effect is moderated by the biomarker
value. Simulation studies are the primary tool for evaluating the
statistical power of proposed trial designs to detect such
interactions. These simulations require a data-generating process
(DGP) that encodes the biomarker-treatment interaction mechanism.

The statistical literature on treatment effect heterogeneity
(Rothwell 2005; Kent et al. 2020) distinguishes between
\textit{qualitative} interactions (the direction of the treatment
effect reverses across biomarker strata) and \textit{quantitative}
interactions (the magnitude varies but the direction is
consistent). Simulation frameworks must choose how to generate
this heterogeneity.

The dominant approach in the clinical trial simulation literature
is direct mean moderation: the biomarker enters the outcome model
as an interaction term that explicitly scales the treatment effect.
This approach is used in essentially all simulation frameworks for
biomarker-stratified designs (Simon 2010), adaptive enrichment
trials (Freidlin \& Korn 2014), basket and umbrella trials (Renfro
\& Sargent 2017), and platform trials. Recent simulation studies
of biomarker-treatment interaction estimation in randomized trials
with prognostic covariates (Haller \& Ulm 2018) likewise adopt
the standard regression specification with explicit interaction
terms. N-of-1 simulation studies (Zucker et al. 1997; Araujo
et al. 2016; Duan et al. 2013) similarly use hierarchical
random-effects models in which the biomarker predicts the
individual treatment effect through the mean structure. Notably,
other methodological work on power and design for N-of-1 trials
with serial correlation and carryover (Wang \& Schork 2019;
Schork 2022) also specifies the treatment effect through a linear
mean structure, indicating that the multivariate-normal
differential-correlation parameterization is not a research-program
convention but a paper-specific design choice in Hendrickson
et al. (2020). To the authors' knowledge, no other clinical trial
simulation framework generates the biomarker-treatment interaction
through differential correlation in a multivariate normal joint
distribution rather than through direct mean moderation.

Both approaches were encountered in the development of the
pmsimstats simulation framework for N-of-1 clinical trials
(Hendrickson et al. 2020). The original publication code uses
the MVN differential correlation approach, an unusual choice
that, as shown here, has substantive consequences for power
estimation under carryover. A comprehensive audit revealed that
a parallel tidyverse reimplementation of the same simulation
had inadvertently adopted the standard direct mean moderation
approach. Both implementations produced plausible power
estimates, but they yielded qualitatively different conclusions
about how carryover effects influence statistical power. Tracing
this discrepancy to its root cause revealed the architectural
difference described in this paper, which has not, to the
authors' knowledge, been explicitly discussed in the simulation
methodology literature.

This paper presents the first systematic comparison of these two
DGP architectures for biomarker-treatment interaction power
estimation under carryover. We formalize the distinction,
demonstrate its consequences with simulation results, and
provide guidance for investigators designing simulation studies
for predictive biomarker-moderated treatment effects.

\textbf{Related documents in this series.} The mathematical
derivation, code-level implementation details, four
implementation attempts (three failed and one correct), and the
empirical evaluation of the original Hendrickson (2020)
parameter grid under Architecture~A are documented in
\texttt{19-mean-moderation-implementation-notes.tex} (with a
brief markdown summary in
\texttt{19-mean-moderation-implementation-notes.md}). A
condensed empirical-results summary suitable for collaborator
presentation is provided in
\texttt{21-architecture-comparison-summary.md}.

\bigskip\noindent\rule{\textwidth}{0.4pt}\bigskip

\section{2. Two DGP Architectures}

\textbf{Notation.} Throughout this paper, $B_i$ denotes the
biomarker value of participant $i$, $D_{it}$ the binary on-drug
indicator at timepoint $t$ (1 on drug, 0 off drug), and
$BR_{it}$ the biological response component of the outcome.
These three symbols describe the data-generating-process layer.
The analysis-model layer uses a related but distinct continuous
drug exposure variable, \texttt{Dbc}, defined in Section 3.1.
Code-style identifiers in typewriter font (\texttt{bm},
\texttt{Dbc}, \texttt{BR}) refer to specific variables in the
pmsimstats codebase or to R formula syntax; mathematical
symbols are used otherwise.

\subsection{2.1 Architecture A: Direct Mean Moderation}

In this approach, the biomarker enters the mean structure of the
outcome directly. The treatment effect for participant $i$ at
timepoint $t$ is:

\[Y_{it} = \mu_0 + \beta_1 \cdot D_{it} \cdot (1 + \beta_{bm}
\cdot B_i) + \epsilon_{it}\]

where $D_{it}$ is the treatment indicator (or a continuous measure
of drug exposure), $B_i$ is the participant's biomarker value
(typically centered), and $\beta_{bm}$ is the biomarker moderation
parameter. The interaction is explicit in the mean: participants
with higher biomarker values have proportionally larger treatment
effects.

Key properties:

\begin{itemize}
\item The interaction is \textbf{deterministic} given the biomarker value
  and treatment status.
\item The biomarker moderation $\beta_{bm}$ is a \textbf{population-level
  parameter} that applies uniformly to all participants.
\item The signal for the interaction is in the \textbf{first moment}
  (conditional mean) of the outcome distribution.
\item Adding noise, random effects, or carryover to $D_{it}$ scales
  the entire treatment effect (including its biomarker-dependent
  component) proportionally. The \textbf{ratio} of drug effect between
  high-biomarker and low-biomarker participants is preserved.
\end{itemize}

The multiplicative form above maps directly onto the standard
additive regression form for predictive biomarker effects used
in the trial methodology literature (e.g., the effect modeling
equation in Kent et al. 2020):
\[Y_{it} = \beta_0 + \beta_1 D_{it} + \beta_2 B_i + \beta_3
D_{it} B_i + \epsilon_{it}\]
The interaction coefficient $\beta_3$ in the additive form
corresponds to the product $\beta_1 \cdot \beta_{bm}$ in the
multiplicative form. The additive form additionally accommodates
a prognostic main effect $\beta_2$ for the biomarker, which the
multiplicative form sets to zero by construction. For the
purpose of power calculations on the interaction coefficient,
the two forms are interchangeable up to this prognostic-effect
distinction, and both fall within Architecture A.

This architecture is used in the pedagogical simulation
(\texttt{implementations/simple/simulation.R}) and in the exploratory carryover
analyses (\texttt{simulation\_carryover\_spectrum.R}) of the
pmsimstats project.

\subsection{2.2 Architecture B: Differential Correlation (MVN)}

In this approach, the biomarker and the biological response (BR)
component are jointly drawn from a multivariate normal distribution
with treatment-state-dependent correlation:

\[\begin{pmatrix} B_i \\ BR_{it} \end{pmatrix} \sim
\text{MVN}\left(\begin{pmatrix} \mu_B \\ \mu_{BR}(t, D_{it})
\end{pmatrix}, \Sigma(D_{it})\right)\]

where the covariance matrix $\Sigma$ has the following
treatment-state-dependent BM-BR correlation:

\[\text{Cor}(B_i, BR_{it}) = \begin{cases}
c_{bm} & \text{if } D_{it} = 1 \text{ (on drug)} \\
c_{bm} \cdot e^{-\lambda \cdot t_{sd}} &
\text{if } D_{it} = 0 \text{ (off drug)}
\end{cases}\]

where $t_{sd}$ is time since drug discontinuation. The
exponential-decay form is the default; the limiting cases
$\lambda = 0$ (no decay; correlation persists indefinitely
during the off-drug phase) and $\lambda \to \infty$ (immediate
decay; correlation drops to zero the instant the drug is
discontinued) recover the two no-carryover model variants.

The interaction is not in the mean structure (the population mean
of $BR_{it}$ is the same for all participants with the same
treatment history). Instead, it emerges from the \textbf{conditional
distribution}: given a participant's biomarker value $B_i = b$,
the conditional expectation of their biological response is:

\[E[BR_{it} | B_i = b, D_{it} = 1] = \mu_{BR}(t) + c_{bm}
\cdot \frac{\sigma_{BR}}{\sigma_B} \cdot (b - \mu_B)\]

This conditional expectation has an interaction-like structure:
the biomarker value $b$ modulates the expected BR. The effective
correlation is $c_{bm}$ on drug and $c_{bm} \cdot
e^{-\lambda t_{sd}}$ during the washout; the moderation therefore
persists with attenuated strength during off-drug periods and
vanishes only in the limit of complete decay.

Key properties:

\begin{itemize}
\item The interaction is \textbf{probabilistic}: it manifests as a
  tendency, not a deterministic scaling.
\item The biomarker moderation $c_{bm}$ is a \textbf{correlation
  parameter} that governs the strength of the association.
\item The signal for the interaction is in the \textbf{second moment}
  (covariance structure) of the joint distribution.
\item Carryover effects in the DGP inflate the BR means during
  off-drug periods, making them resemble on-drug means. This
  reduces the \textbf{differential} in the correlation structure
  between treatment states, weakening the signal that the
  analysis model detects.
\end{itemize}

This architecture is used in the Hendrickson et al. (2020)
publication code and the audited \texttt{R/} package in pmsimstats-ng.

\subsection{2.3 Mathematical Comparison}

Under Architecture A, the expected outcome difference between
two participants with biomarker values $b_1$ and $b_2$ at the
same drug exposure level $D$ is:

\[E[Y \mid b_1, D] - E[Y \mid b_2, D] = \beta_1 \cdot \beta_{bm}
\cdot (b_1 - b_2) \cdot D\]

The absolute magnitude of this difference scales linearly with
$D$ and therefore shrinks during off-drug periods when $D$ is
replaced by an exposure-decayed value $D \cdot \text{carryover}$.
However, the \textbf{ratio} of the drug effect between
high-biomarker and low-biomarker participants is preserved at
every level of exposure. The biomarker-by-treatment signal
contracts in amplitude but does not lose its identifiability
under carryover.

Under Architecture B, the analogous quantity is:

\[E[BR | b_1, D=1] - E[BR | b_2, D=1] = c_{bm} \cdot
\frac{\sigma_{BR}}{\sigma_B} \cdot (b_1 - b_2)\]

During off-drug periods with carryover:

\[E[BR | b, D=0, t_{sd}] = \mu_{BR,\text{off}}(t_{sd}) +
c_{bm} \cdot e^{-\lambda t_{sd}} \cdot
\frac{\sigma_{BR}}{\sigma_B} \cdot (b - \mu_B)\]

The biomarker-dependent component decays exponentially with
time off drug. As $t_{sd}$ increases, the off-drug conditional
expectation converges to $\mu_{BR,\text{off}}$ for all
biomarker values, eliminating the differential signal.

\bigskip\noindent\rule{\textwidth}{0.4pt}\bigskip

\section{3. Consequences for Statistical Power Under Carryover}

\subsection{3.1 Empirical demonstration}

Identical simulation configurations were run under both
architectures using the pmsimstats-ng framework. The DGP
parameters matched as closely as possible between the two
approaches. Critically, the non-biologic response components
(time-varying response, placebo-belief accumulation, and
residual noise) are generated from the same multivariate normal
draw with identical means, variances, within-factor AR(1)
autocorrelation, and cross-factor correlations in both
architectures. The two DGPs differ \emph{only} in how the
biomarker-to-biological-response (BM-BR) linkage is encoded:
Architecture~B embeds it in the off-diagonal $c_{bm}$ entries of
$\Sigma$, whereas Architecture~A omits those entries and adds a
post-draw shift $c_{bm} \cdot B_i^{*} \cdot \sigma_{BR}$ to
on-drug BR columns (see
\texttt{19-mean-moderation-implementation-notes.tex} for the
scaling derivation). The divergent carryover behavior reported
below is therefore attributable to the BM-BR parameterization,
not to any difference in how the placebo, time-varying, or
residual-noise components are generated. The analysis model was
the same in both cases: an \texttt{nlme::lme} fit with a
\texttt{corCAR1} continuous-time autocorrelation structure and a
single biomarker-by-drug interaction term, written
\texttt{bm:Dbc} in R formula syntax.

The analysis-model drug exposure variable \texttt{Dbc} is the
\emph{continuous} analogue of the binary DGP indicator $D_{it}$.
It equals $1$ during on-drug periods and $\exp(-\lambda t_{sd})$
during off-drug periods, where $t_{sd}$ is time since drug
discontinuation and $\lambda$ is the analysis-model carryover
decay rate (typically derived from the drug's pharmacokinetic
half-life as $\lambda = \ln 2 / t_{1/2}$, in which case it
matches the DGP decay rate of Section 2.2 by construction).
Whereas $D_{it}$ flips abruptly between 0 and 1, \texttt{Dbc}
captures the smooth shrinkage of residual drug effect during
washout. The interaction coefficient on \texttt{bm:Dbc} is
therefore the test of whether the biomarker moderates the
exposure-decayed drug effect.

\textbf{Setup:}

\begin{itemize}
\item Trial designs: OL, CO, N-of-1 (Hybrid), OL+BDC
\item Sample size: $N = 70$
\item Biomarker correlation / moderation: $c_{bm} = 0.45$ /
  $\beta_{bm} = 0.45$
\item DGP carryover half-life: $t_{1/2} \in \{0, 0.5, 1.0\}$ weeks
\item Analysis: matched \texttt{Dbc} with same half-life
\item Replicates: 50 per cell
\end{itemize}

The 50-replicate setting is intentionally modest and the tables
below should be read as a proof of concept rather than precise
power estimates. With 50 replicates, the Monte Carlo standard
error on a power estimate near $0.5$ is approximately
$\sqrt{0.5 \cdot 0.5 / 50} \approx 0.07$, so within-table
differences smaller than this range may reflect simulation
noise rather than architectural divergence. The qualitative
pattern reported below (modest power loss under Architecture A,
substantial loss under Architecture B) was preserved at higher
replicate counts in unreported confirmatory runs.

\textbf{Results under Architecture A (direct mean moderation):}

\begin{table}[htbp]
\centering
\begin{tabular}{lccc}
\toprule
Design & $t_{1/2} = 0$ & $t_{1/2} = 0.5$ & $t_{1/2} = 1.0$ \\
\midrule
CO & 0.38 & 0.36 & 0.34 \\
N-of-1 & 0.74 & 0.72 & 0.68 \\
OL+BDC & 0.62 & 0.58 & 0.54 \\
\bottomrule
\end{tabular}
\end{table}

Power declines modestly (\textasciitilde{}10-15\% relative) with increasing
carryover.

\textbf{Results under Architecture B (differential correlation, MVN):}

\begin{table}[htbp]
\centering
\begin{tabular}{lccc}
\toprule
Design & $t_{1/2} = 0$ & $t_{1/2} = 0.5$ & $t_{1/2} = 1.0$ \\
\midrule
CO & 0.40 & 0.40 & 0.32 \\
N-of-1 & 0.82 & 0.64 & 0.50 \\
OL+BDC & 0.62 & 0.38 & 0.24 \\
\bottomrule
\end{tabular}
\end{table}

Power declines substantially (\textasciitilde{}40-60\% relative) with increasing
carryover.

\subsection{3.2 Why the architectures diverge}

\textbf{Under Architecture A}, carryover reduces the magnitude
of the drug exposure variable seen by the analysis model
(\texttt{Dbc} drops below 1 during washout), and the mean BR at
off-drug timepoints is inflated by the residual drug effect
exactly as in Architecture B. Both mechanisms compress the
between-state treatment contrast and are shared by the two
architectures. What Architecture A lacks is the second channel
of signal loss: the biomarker moderation coefficient
$\beta_{bm}$ is a fixed population parameter, the biomarker
enters the mean structure directly rather than through a
correlation, and the noise structure is independent of drug
state. The biomarker-by-drug interaction term therefore has
reduced variance (because \texttt{Dbc} has less contrast between
drug states), but the signal-to-noise ratio degrades only
modestly.

\textbf{Under Architecture B}, carryover operates on two levels
simultaneously:

\begin{enumerate}
\item \textbf{Mean blurring}: The BR means during off-drug periods are
  inflated by residual carryover, making them resemble on-drug
  BR means. This reduces the treatment contrast that the analysis
  model uses to distinguish drug effect from no-drug. This channel
  is shared with Architecture A.

\item \textbf{Correlation erosion}: The BM-BR correlation during off-drug
  periods decays exponentially with $t_{sd}$. This is not just a
  statistical modeling choice -- it is a structural consequence
  of the DGP. The correlation between biomarker and BR reflects
  the degree to which the drug effect is active; as the drug
  washes out, the correlation weakens because the drug-mediated
  component of BR variance shrinks relative to the
  drug-independent component. This channel is \emph{unique} to
  Architecture B, because only here does the biomarker signal
  live in the second moment of the joint distribution. Under
  Architecture A the drug-mediated and drug-independent
  components of BR variance move in tandem -- both are scaled by
  the same \texttt{Dbc}-weighted drug effect -- so their ratio,
  and hence the identifiability of the interaction, is preserved
  even as amplitudes shrink.
\end{enumerate}

The analysis model's interaction coefficient (the coefficient
on \texttt{bm:Dbc}) depends on the covariance between $B_i$ and
\texttt{Dbc}-weighted outcomes. Under Architecture B, this
covariance is being attacked from both sides: the treatment
contrast in \texttt{Dbc} is compressed, and the BM-BR
correlation that generates the covariance signal is
simultaneously decaying.

\bigskip\noindent\rule{\textwidth}{0.4pt}\bigskip

\section{4. Biological Assumptions}

The choice between architectures is not merely a statistical
convenience -- it reflects different assumptions about the
biological mechanism by which the biomarker moderates treatment
response.

\subsection{4.1 When Architecture A is appropriate}

Architecture A (direct mean moderation) is appropriate when the
biomarker governs the \textbf{magnitude of the drug's biological
effect} for each individual participant. The drug effect for
participant $i$ is scaled by their biomarker value:
$\text{effect}_i = f(B_i) \cdot \text{drug}$. The scaling is
deterministic in the population-mean sense --- any two
participants sharing the same $B_i$ have the same expected
response --- but individual outcomes still carry the usual noise,
random-effect, and residual-biological variability. The label
``deterministic'' refers to the mean structure of the DGP, not to
the realized response of any single participant.

\textbf{Clinical examples:}

\begin{itemize}
\item \textbf{Renal clearance of a renally-excreted drug.} Estimated
  glomerular filtration rate (eGFR) directly governs the
  elimination rate of drugs such as digoxin or aminoglycoside
  antibiotics. Participants with higher eGFR clear the drug more
  rapidly, achieving lower steady-state concentrations and a
  proportionally smaller drug effect. The biomarker-drug
  relationship is mechanistic and deterministic: given the
  biomarker value, the effective dose at steady state is fully
  determined (up to random noise).

\item \textbf{Receptor density moderation.} A biomarker that measures
  target receptor density (e.g., PET imaging of receptor
  availability for an antagonist) directly determines the
  pharmacodynamic response to that antagonist. More receptors
  means more drug effect, proportionally.

\item \textbf{Genetic metabolizer status.} CYP2D6 metabolizer phenotype
  determines the rate of drug activation or clearance. Poor
  metabolizers of a prodrug receive less active compound,
  reducing their treatment effect proportionally.

\item \textbf{Dose-response with biomarker-dependent effective dose.}
  When the biomarker determines the effective dose (through
  body composition, protein binding, or volume of distribution),
  the treatment effect scales with the biomarker through the
  dose-response curve.
\end{itemize}

In all these cases, the biomarker's moderating effect is
preserved under carryover because the residual drug effect during
off-drug periods maintains the same proportional relationship
with the biomarker. If participant A has twice the drug effect of
participant B when on drug, they will also have approximately
twice the residual effect during the washout period.

\subsection{4.2 When Architecture B is appropriate}

Architecture B (differential correlation) is appropriate when
the biomarker predicts \textbf{which participants are likely to
respond}, but the magnitude of any individual's response is not
deterministically governed by the biomarker. The biomarker and
the drug response are associated at the population level, but
the association is mediated by latent factors (disease subtype,
neurobiological heterogeneity, comorbidities) that the biomarker
imperfectly indexes.

\textbf{Clinical examples:}

\begin{itemize}
\item \textbf{Blood pressure as a PTSD subtype marker.} Elevated resting
  blood pressure may index a noradrenergic-predominant PTSD
  phenotype (Hendrickson et al. 2020) that is more likely to
  respond to prazosin (an alpha-1 adrenergic antagonist). The
  blood pressure does not \textit{cause} the drug to work better --
  it correlates with an underlying neurobiological state that
  determines drug responsiveness. Two participants with
  identical blood pressure may have very different responses
  because blood pressure is an imperfect proxy for the
  noradrenergic state.

\item \textbf{Baseline disease severity as a treatment predictor.} Trials
  in Alzheimer's disease, depression, and chronic pain often
  find that patients with more severe baseline symptoms show
  larger treatment effects (the `floor effect' or `regression
  to the mean' phenomenon). The baseline severity score
  correlates with treatment response, but the relationship is
  statistical (population-level tendency), not deterministic.

\item \textbf{Inflammatory biomarkers in psychiatry.} C-reactive protein
  (CRP) or interleukin-6 levels may predict response to
  anti-inflammatory augmentation of antidepressants (Raison
  et al. 2013). High-CRP patients are more likely to respond,
  but the relationship is probabilistic -- many high-CRP
  patients do not respond, and some low-CRP patients do.

\item \textbf{Genetic risk scores.} Polygenic risk scores predict disease
  trajectory and treatment response, but explain only a fraction
  of the variance. The remainder reflects unmeasured genetic,
  epigenetic, and environmental factors.
\end{itemize}

In these cases, carryover has a different implication: as the
drug washes out, the correlation between biomarker and response
weakens because the drug-mediated component of response variance
diminishes. The biomarker's predictive power is inherently tied
to the presence of active drug effect. Without drug, the
biomarker is just a biomarker -- it does not predict the
non-drug component of symptom change.

The latent-subtype framing raises a legitimate modeling
question: if the biomarker is really a noisy indicator of an
underlying responder class, the biologically faithful DGP would
be an explicit finite mixture in which each participant is drawn
from one of two (or more) response distributions with
class-membership probability a function of $B_i$. Under such a
mixture DGP the individual-level moderation is discrete, not
probabilistic, and the MVN differential-correlation
parameterization is best understood as a tractable
second-moment approximation to a richer mixture model rather
than as the generative mechanism itself. The practical
consequence is that Architecture~B as implemented here captures
the population-level covariance implications of a latent-class
mechanism without committing to its full generative form; this
is adequate for power calculations that test for the presence
of an interaction but should not be read as an endorsement of
MVN differential correlation as the correct mechanistic model.
Whether the observed power loss under carryover is better
mitigated by analysis strategies (Section 6) or by explicit
mixture modeling is an open question.

The latent-subtype framing of Architecture B should also be
adopted with appropriate epistemic caution. Senn (2016) has
argued that
the apparent ``personal element in response to treatment'' is
often a variance-component artifact rather than a genuine
biomarker-by-treatment interaction, and that even modest
within-subject variability can produce convincing-looking but
spurious differential response. Architecture B presupposes that
the latent subtype indexed by the biomarker is real,
identifiable, and stable across the trial period. When that
presupposition is shaky, the apparent power loss under carryover
documented in Section 3 may overstate what could be recovered
under any analysis model, because the underlying signal it is
attempting to estimate may itself be smaller than assumed.

\subsection{4.3 The prazosin-PTSD case}

The motivating application for pmsimstats is the active-duty
prazosin trial by Raskind et al. (2013), whose SBP-stratified
secondary analysis is the subject of ongoing reanalysis by the
pmsimstats team.
The biomarker is resting systolic blood pressure (SBP) and the
outcome is CAPS (Clinician-Administered PTSD Scale) score
change.

The clinical evidence base for prazosin in PTSD is mixed. The
earlier single-site trials by Raskind and colleagues (2003,
2007, 2013) reported benefit on nightmares and global PTSD
symptoms in veteran and active-duty samples. The subsequent
multisite Veterans Affairs Cooperative Study Program trial
(Raskind et al. 2018), known as PACT, did not replicate the
effect: prazosin failed to outperform placebo on the primary
endpoints. The negative PACT result has been interpreted by
several investigators as consistent with effect heterogeneity --
prazosin works in some patients but not others -- which
motivates the search for a predictive biomarker that would
identify the responder subgroup. Resting SBP, as a noninvasive
proxy for central noradrenergic tone, is the leading candidate.

The biological hypothesis is that elevated SBP indexes
heightened central noradrenergic tone, which characterizes a
subtype of PTSD that is preferentially responsive to alpha-1
adrenergic blockade by prazosin. This is an Architecture B
scenario: SBP is a proxy for an underlying neurobiological
state, not a direct determinant of drug pharmacokinetics. Two
patients with the same SBP may differ in their actual
noradrenergic tone, alpha-1 receptor distribution, and
downstream signaling, leading to different treatment responses
despite identical biomarker values.

The Hendrickson et al. (2020) DGP uses Architecture B (MVN
differential correlation) for this clinical context. The
finding that carryover substantially reduces power under this
architecture is therefore clinically relevant: it suggests that
trial designs with longer off-drug periods will have
meaningfully less power to detect the SBP-prazosin interaction
than designs with shorter off-drug periods, and this effect is
larger than what a direct mean moderation simulation would
predict.

\bigskip\noindent\rule{\textwidth}{0.4pt}\bigskip

\section{5. Implications for Simulation Study Design}

\subsection{5.1 Choosing the appropriate architecture}

The choice of DGP architecture should be guided by the biological
mechanism hypothesized for the biomarker-treatment interaction:

\begin{table}[htbp]
\centering
\begin{tabular}{lcc}
\toprule
Criterion & Architecture A & Architecture B \\
\midrule
Biomarker role & Causal mediator & Statistical predictor \\
Mechanism & Deterministic scaling & Probabilistic association \\
Signal location & Mean structure & Covariance structure \\
Carryover impact on power & Modest (\textasciitilde{}10-15\%) & Substantial (\textasciitilde{}40-60\%) \\
Appropriate when & \parbox[t]{3.5cm}{Biomarker determines\\effective dose or PK} & \parbox[t]{3.5cm}{Biomarker indexes a\\latent subtype} \\
\bottomrule
\end{tabular}
\end{table}

\subsection{5.2 Reporting requirements}

Simulation studies evaluating biomarker-moderated treatment effects
should explicitly state:

\begin{enumerate}
\item Whether the biomarker-treatment interaction is generated through
  mean moderation or differential correlation.
\item The biological rationale for the chosen mechanism.
\item Whether the carryover model (if present) operates on the
  interaction signal consistently with the chosen architecture.
\item Sensitivity analyses under the alternative architecture, if the
  biological mechanism is uncertain.
\end{enumerate}

\subsection{5.3 When the choice is uncertain}

If the biological mechanism is ambiguous (as it often is in early
biomarker development), investigators should:

\begin{enumerate}
\item Run the simulation under both architectures and report the
  range of power estimates.
\item Architecture B produces more conservative (lower)
  power estimates under carryover, making it the safer default
  for trial design decisions.
\item Design the trial to minimize carryover (adequate washout
  periods) to reduce sensitivity to the DGP architecture choice.
\end{enumerate}

\bigskip\noindent\rule{\textwidth}{0.4pt}\bigskip

\section{6. Alternative Analysis Strategies Under Architecture B}

The power loss documented in Section 3 arises from two distinct
sources: (1) compressed variance of the treatment indicator
\texttt{Dbc} when carryover is present, and (2) erosion of the
BM-BR correlation during off-drug periods. Source 1 is a
property of
the analysis model parameterization and can potentially be
addressed by choosing a different predictor. Source 2 is a
property of the data itself -- the signal is genuinely weaker in
off-drug observations -- and no analysis model can recover signal
that is not present. However, analysis strategies that are
selective about \textit{which observations contribute to the test} or
\textit{how they are weighted} can mitigate the damage.

\subsection{6.1 Restrict analysis to on-drug observations}

The most direct approach is to discard off-drug timepoints
entirely and test the biomarker-treatment interaction using only
on-drug observations. Every retained observation has the full
BM-BR correlation ($c_{bm}$), so correlation erosion is
eliminated. The costs are twofold: a reduction in effective
sample size (fewer observations per participant) and, more
importantly, the loss of the off-drug reference that separates
the biological response from placebo and time-varying
background response. Without off-drug observations the
biological, placebo, and time-varying components are confounded,
which is why the Open-Label design -- despite having the most
on-drug timepoints -- does not dominate the other designs in
the Section 3 results. For designs with many
on-drug timepoints (OL has 8, the Hybrid design has 5-6 per
path), the net effect may be positive: the per-observation
signal-to-noise ratio is maximal, and the analysis model
simplifies to \texttt{Sx \textasciitilde{} bm + t + bm:t} since there is no treatment
contrast to model. The interaction is detected through the
correlation between biomarker and the time-on-drug trajectory.

This approach is most promising for the Open-Label design (which
has no off-drug observations regardless) and the Hybrid design
(which has a long initial on-drug phase). It is least useful for
the Crossover design, where discarding half the observations
substantially reduces statistical power.

\subsection{6.2 Weighted analysis}

Rather than discarding off-drug observations entirely, a weighted
analysis retains all observations but weights them by their
estimated drug exposure level. On-drug observations receive
weight 1. Off-drug observations receive weight proportional to
the expected residual drug effect:

\[w_t = e^{-\lambda \cdot t_{sd}}\]

This downweights observations where the BM-BR correlation has
decayed (and where they contribute more noise than signal to the
interaction estimate) without discarding them entirely. The
rationale is information-theoretic: observations with higher drug
exposure carry more information about the biomarker-treatment
interaction because the correlation signal is stronger. This is
straightforward to implement in \texttt{nlme::lme} using the \texttt{weights}
argument.

\subsection{6.3 Within-subject contrast}

For designs with both on-drug and off-drug periods (CO, Hybrid,
OL+BDC), a summary-measure approach computes a per-participant
contrast: the difference between mean on-drug response and mean
off-drug response. The biomarker-treatment interaction is then
tested by a simple regression of these contrasts on the biomarker:

\[\bar{Y}_{i,\text{on}} - \bar{Y}_{i,\text{off}} = \alpha +
\gamma \cdot B_i + \epsilon_i\]

This sidesteps the \texttt{Dbc} parameterization and the
repeated-measures model entirely. Carryover affects the
magnitude
of the contrast (making off-drug means closer to on-drug means),
but the \textit{correlation between the contrast and the biomarker} may
be more robust because within-subject averaging reduces noise
before the biomarker association is tested. The approach trades
the efficiency of a full longitudinal model for robustness to
misspecification of the carryover and correlation structure.

\subsection{6.4 Two-stage random slopes}

A two-stage approach first estimates participant-specific
treatment effects, then tests their association with the
biomarker. Stage 1 fits a random-slopes model:

\[Y_{it} = \beta_0 + \beta_1 t + \beta_2 D_{it} + u_{0i} +
u_{1i} D_{it} + \epsilon_{it}\]

where $\beta_2$ is the population mean drug effect and $u_{1i}$
is participant $i$'s random deviation from that mean. The
participant-specific drug effect is then $\beta_2 + u_{1i}$.
Stage 2 regresses the estimated random deviations (or
equivalently, the BLUPs of the participant-specific effects) on
the biomarker:

\[\hat{u}_{1i} = \delta_0 + \delta_1 B_i + \eta_i\]

This separates the within-subject treatment effect estimation
(which uses the longitudinal data and is affected by carryover)
from the between-subject biomarker association (which uses the
cross-sectional relationship between estimated effects and
biomarker values). The approach may be more robust to carryover
because the random slope $u_{1i}$ integrates over the
participant's entire trajectory, including both on-drug and
off-drug phases.

\subsection{6.5 Exclusion of contaminated observations}

The first 1-2 off-drug timepoints carry the most carryover
contamination: the residual drug effect is highest, and the
BM-BR correlation is in the process of decaying. Later off-drug
timepoints, where the drug has fully washed out, provide a
cleaner contrast with the on-drug phase. Excluding the early
off-drug observations and retaining only the later ones yields
a comparison between fully-on and fully-off states, avoiding the
intermediate zone where the analysis model is most strained.

\subsection{6.6 Design-level solutions}

Rather than adapting the analysis model to accommodate carryover,
the trial design itself can be modified to reduce carryover
exposure:

\begin{itemize}
\item \textbf{Longer washout periods} between drug phases eliminate
  residual drug effect before off-drug measurements begin.
\item \textbf{More on-drug timepoints} relative to off-drug increase the
  proportion of high-signal observations.
\item \textbf{Front-loaded drug exposure} (as in the OL+BDC design)
  places a long initial on-drug block before the off-drug
  baseline-during-continued-treatment phase, so that the BM-BR
  correlation has been fully established before any off-drug
  observations are collected. The alternative ordering
  (off-drug first, then on-drug) has the appealing property of
  placing the placebo and time-varying estimation in an
  uncontaminated window but sacrifices on-drug exposure duration
  and forces the BM-BR signal to be re-established after a
  wash-in; the net power trade-off depends on the ratio of
  carryover half-life to total trial length and is explored
  empirically in the Section 3 results.
\end{itemize}

Design-level solutions address the root cause (carryover in the
data) rather than the symptom (power loss in the analysis). Their
cost is typically a longer trial duration or reduced ability to
estimate the carryover dynamics themselves.

\subsection{6.7 Comparative evaluation}

The strategies described above differ in their assumptions,
implementation complexity, and expected power recovery:

\begin{table}[htbp]
\centering
\begin{tabular}{lccp{5cm}}
\toprule
Strategy & Discards data & Complexity & Expected benefit \\
\midrule
On-drug only & Yes & Low & High for OL/Hybrid; poor for CO \\
Weighted & No & Low & Moderate; depends on weight calibration \\
Within-subject contrast & Aggregates & Low & Moderate; robust to misspecification \\
Two-stage random slopes & No & Moderate & Unknown; separates estimation stages \\
Exclude early off-drug & Some & Low & Moderate; improves contrast clarity \\
Design modification & N/A & N/A & High; addresses root cause \\
\bottomrule
\end{tabular}
\end{table}

A systematic simulation comparison of these strategies under
Architecture B, analogous to the factorial comparison of analysis
models in Section 3, would identify which approach best recovers
the power lost to carryover for each trial design. This is a
direction for future work.

\bigskip\noindent\rule{\textwidth}{0.4pt}\bigskip

\section{7. Relationship to the Broader Literature}

\subsection{7.1 Treatment effect heterogeneity}

The statistical literature on treatment effect heterogeneity
(Rothwell 2005; Varadhan et al. 2013) distinguishes between
\textit{predictive} biomarkers (those that modify the treatment
effect) and \textit{prognostic} biomarkers (those that predict
outcome regardless of treatment). The current consensus
statement on predictive HTE analysis is the Predictive
Approaches to Treatment effect Heterogeneity (PATH) Statement
(Kent et al. 2020), which distinguishes two analytic strategies
for identifying predictive biomarker effects in randomized
trials: a \textit{risk modeling} approach in which a multivariable
risk model (developed without a treatment term) is used to
stratify patients and examine variation in absolute treatment
benefit across risk strata, and an \textit{effect modeling} approach
in which the regression includes explicit treatment-by-covariate
interaction terms. PATH expresses caution about data-driven
effect modeling on the grounds of low statistical power,
multiplicity, and overfitting, and favors risk modeling when
the application criteria are met.

The PATH risk-versus-effect-modeling distinction is an
analysis-model axis. The Architecture A versus Architecture B
distinction introduced in this paper is a data-generating-process
axis, and the two are orthogonal. Both PATH analytic strategies
implicitly assume that the underlying treatment effect
heterogeneity is encoded in the mean structure of the outcome
(the heterogeneity is identifiable from differences in
conditional means), which corresponds to Architecture A as a
DGP. Architecture B as a DGP is a separate phenomenon that does
not map cleanly onto either of PATH's analysis-model categories:
the interaction signal resides in the covariance structure
rather than the mean, and the standard regression-based
predictive HTE analysis methods endorsed by PATH may have
reduced power under Architecture B precisely because they were
developed under an implicit Architecture A assumption.
Architecture A can accordingly be characterized as a
\textit{parametric interaction model} (in PATH's effect modeling
sense). Architecture B can be characterized as a \textit{latent
class} or \textit{mixture model} perspective in which the biomarker is
a noisy indicator of class membership.

\subsection{7.2 Prevalence of each architecture}

A survey of the simulation methodology literature reveals that
Architecture A (direct mean moderation) is the near-universal
standard. All simulation frameworks for biomarker-stratified
designs (Simon 2010), adaptive enrichment trials (Freidlin \&
Korn 2014), basket and umbrella trials (Renfro \& Sargent 2017),
and platform trials use direct mean moderation with explicit
interaction terms. Recent simulation work on biomarker-by-
treatment interaction estimation in randomized trials with
prognostic covariates (Haller \& Ulm 2018) likewise adopts the
standard regression specification, focusing on power and bias
of the interaction estimator under different covariate
adjustment strategies. N-of-1 simulation studies (Zucker et al.
1997; Araujo et al. 2016; Duan et al. 2013) use hierarchical
random-effects models in which individual treatment effects are
drawn from a population distribution -- a variant of
Architecture A where the biomarker predicts the random
treatment effect through the mean structure. Methodological
work on power and design for N-of-1 trials with serial
correlation and carryover (Wang \& Schork 2019; Schork 2022)
adopts the same Architecture A specification, even though it
shares senior authorship with Hendrickson et al. (2020). This
suggests that the differential-correlation parameterization in
the latter is a paper-specific design choice rather than a
methodological convention of any laboratory or research program.

Architecture B (differential correlation via MVN) appears
primarily in latent variable and structural equation modeling
contexts (Rizopoulos 2012), Bayesian joint modeling frameworks,
and copula-based simulation studies. In the clinical trial
simulation literature specifically, the Hendrickson et al.
(2020) framework appears to be unique in using differential
correlation as the mechanism for biomarker-treatment
interaction in an N-of-1 trial context.

This predominance of Architecture A has an important implication:
the power estimates reported in the vast majority of
biomarker-moderated trial simulation studies may be optimistic
for biomarkers that operate through an Architecture B mechanism,
because these studies do not account for the additional power
loss that carryover produces when the interaction signal resides
in the covariance structure rather than the mean structure.

\subsection{7.3 Crossover and N-of-1 trial methodology}

The crossover trial literature (Senn 2002; Jones \& Kenward
2014) addresses carryover extensively but does not distinguish
between architectures for biomarker-treatment interaction
modeling. More recent methodological work on testing for
carryover effects via mixed-effects models (Sturdevant \& Lumley
2021) likewise treats carryover as a nuisance parameter to be
estimated and adjusted for, without engaging with how carryover
interacts with biomarker-by-treatment effect signals. The
standard recommendation -- always include carryover in the
analysis model or use adequate washout -- applies to both
architectures but has different power implications under each.
The finding that Architecture B is more sensitive to carryover
suggests that crossover designs with short washout periods may
be less suitable for detecting correlation-based biomarker
interactions than previously recognized.

\subsection{7.4 Precision medicine trial design}

Enrichment and biomarker-stratified designs (Simon 2010;
Freidlin \& Korn 2014) use Architecture A exclusively for power
calculations. Three families of design are particularly worth
considering in light of the architectural distinction.

\textbf{Adaptive enrichment trials} adjust the eligible
population mid-trial based on emerging evidence of a
biomarker-treatment interaction. The interim power calculations
that drive these adaptive decisions are universally framed in
the Architecture A regression idiom. If the true biomarker
mechanism is closer to Architecture B and any participants pass
through washout intervals between dose-finding stages, the
realized power for the interaction test will be lower than the
adaptation rule expects, biasing decisions toward premature
futility stopping or toward overly narrow enrichment criteria.
The discrepancy is most acute when the biomarker mechanism is
uncertain at the design stage, which is precisely the situation
in which adaptive enrichment is most attractive.

\textbf{Treatment-switching designs} (most crossover, basket,
and platform trials, as well as N-of-1 hybrid designs) cycle
individual participants through different drug states. Under
Architecture B, each washout phase erodes the BM-BR correlation,
so the biomarker-treatment interaction signal in the second and
subsequent treatment periods is structurally weaker than in the
first. The power loss documented in Section 3 applies here
directly. By contrast, parallel-group enrichment designs --- in
which each participant remains on a single treatment for the
duration of the trial --- escape this loss because there is no
off-drug phase during which the correlation can decay; for
these designs, Architectures A and B yield similar power
estimates. The architectural choice therefore matters most in
exactly the designs where the biomarker mechanism is most
likely to be of clinical interest, namely those that
deliberately expose each participant to both treatment states
in order to estimate within-subject responsiveness.

\textbf{Basket and umbrella trials} (Renfro \& Sargent 2017)
introduce an additional layer because their power calculations
average across multiple biomarker-defined arms. If even a
subset of the arms operate through an Architecture B
mechanism while the trial-wide power calculation assumes
Architecture A throughout, the aggregated power estimates will
be biased optimistically in a way that is difficult to detect
from trial-level results alone. Subgroup-specific power
diagnostics under both architectures should accompany the
primary planning analysis whenever the underlying biology is
heterogeneous across baskets.

In all three settings the practical recommendation parallels
that of Section 5.3: when the biological mechanism is uncertain,
run the power calculation under both architectures and treat
the Architecture B estimate as the more conservative anchor for
sample size planning. This recommendation is consistent with
the PATH Statement (Kent et al. 2020), which itself counsels
caution about data-driven discovery of treatment-by-covariate
interactions and prefers risk modeling approaches when the
application criteria are met.

\bigskip\noindent\rule{\textwidth}{0.4pt}\bigskip

\section{8. Conclusions}

\begin{enumerate}
\item The choice between direct mean moderation (Architecture A) and
  differential correlation (Architecture B) for generating
  biomarker-treatment interactions in clinical trial simulations
  has substantive consequences for power estimates under
  carryover.

\item Under Architecture A, carryover reduces power modestly
  (\textasciitilde{}10-15\%) because the proportional biomarker-drug relationship
  is preserved during washout. Under Architecture B, carryover
  reduces power substantially (\textasciitilde{}40-60\%) because the differential
  correlation signal erodes as the drug effect wanes.

\item The choice should be guided by the hypothesized biological
  mechanism: Architecture A for biomarkers that directly
  determine drug pharmacokinetics or pharmacodynamics;
  Architecture B for biomarkers that statistically predict
  treatment response through latent subtype membership.

\item For the prazosin-PTSD application with blood pressure as the
  predictive biomarker, Architecture B (MVN differential
  correlation) is the more appropriate model, and the resulting
  power sensitivity to carryover should inform trial design
  decisions regarding off-drug period duration.

\item The existing clinical trial simulation literature uses
  Architecture A almost exclusively. The Hendrickson et al.
  (2020) MVN approach is, to the authors' knowledge, unique in the
  N-of-1 trial context. This means that published power
  estimates for biomarker-moderated crossover and N-of-1
  designs may be systematically optimistic for biomarkers
  that operate through a correlation-based (Architecture B)
  mechanism, because they do not account for the carryover
  sensitivity documented here.

\item Simulation studies for biomarker-moderated treatment effects
  should explicitly report their DGP architecture and consider
  sensitivity analyses under the alternative when the biological
  mechanism is uncertain.
\end{enumerate}

\bigskip\noindent\rule{\textwidth}{0.4pt}\bigskip

\section{References}

\sloppy

Araujo A, Julious S, Senn S. Understanding variation in sets of
N-of-1 trials. \textit{PLoS ONE}. 2016;11(12):e0167167.

Duan N, Kravitz RL, Schmid CH. Single-patient (N-of-1) trials: a
pragmatic clinical decision methodology. \textit{J Clin Epidemiol}.
2013;66(8):S21-S28.

Freidlin B, Korn EL. Biomarker enrichment strategies: matching
trial design to biomarker credentials. \textit{Nat Rev Clin Oncol}.
2014;11(2):81-90.

Haller B, Ulm K. A simulation study on estimating
biomarker-treatment interaction effects in randomized trials
with prognostic variables. \textit{Trials}. 2018;19(1):128.

Hendrickson RC, Thomas RG, Schork NJ, Raskind MA. Optimizing
aggregated N-of-1 trial designs for predictive biomarker
validation: statistical methods and practical considerations.
\textit{Front Digit Health}. 2020;2:13.

Jones B, Kenward MG. \textit{Design and Analysis of Cross-Over Trials}.
3rd ed. Boca Raton, FL: CRC Press; 2014.

Kent DM, Paulus JK, van Klaveren D, et al. The Predictive
Approaches to Treatment effect Heterogeneity (PATH) Statement.
\textit{Ann Intern Med}. 2020;172(1):35-45.

Raison CL, Rutherford RE, Woolwine BJ, et al. A randomized
controlled trial of the tumor necrosis factor antagonist
infliximab for treatment-resistant depression: the role of
baseline inflammatory biomarkers. \textit{JAMA Psychiatry}.
2013;70(1):31-41.

Raskind MA, Peskind ER, Kanter ED, et al. Reduction of
nightmares and other PTSD symptoms in combat veterans by
prazosin: a placebo-controlled study. \textit{Am J Psychiatry}.
2003;160(2):371-373.

Raskind MA, Peskind ER, Hoff DJ, et al. A parallel group
placebo-controlled study of prazosin for trauma nightmares and
sleep disturbance in combat veterans with post-traumatic stress
disorder. \textit{Biol Psychiatry}. 2007;61(8):928-934.

Raskind MA, Peterson K, Williams T, et al. A trial of prazosin
for combat trauma PTSD with nightmares in active-duty soldiers
returned from Iraq and Afghanistan. \textit{Am J Psychiatry}.
2013;170(9):1003-1010.

Raskind MA, Peskind ER, Chow B, et al. Trial of prazosin for
post-traumatic stress disorder in military veterans.
\textit{N Engl J Med}. 2018;378(6):507-517.

Renfro LA, Sargent DJ. Statistical controversies in clinical
research: basket trials, umbrella trials, and other master
protocols. \textit{Ann Oncol}. 2017;28(1):34-43.

Rizopoulos D. \textit{Joint Models for Longitudinal and
Time-to-Event Data: With Applications in R}. Boca Raton, FL:
CRC Press; 2012.

Rothwell PM. Treating individuals 5. Subgroup analysis in
randomised controlled trials: importance, indications, and
interpretation. \textit{Lancet}. 2005;365(9454):176-186.

Schork NJ. Accommodating serial correlation and sequential
design elements in personalized studies and aggregated
personalized studies. \textit{Harv Data Sci Rev}. 2022;Special
Issue 3.

Senn S. \textit{Cross-over Trials in Clinical Research}. 2nd ed.
Chichester, UK: John Wiley \& Sons; 2002.

Senn S. Mastering variation: variance components and
personalised medicine. \textit{Stat Med}. 2016;35(7):966-977.

Simon R. Clinical trial designs for evaluating the medical
utility of prognostic and predictive biomarkers in oncology.
\textit{Per Med}. 2010;7(1):33-47.

Sturdevant SG, Lumley T. Statistical methods for testing
carryover effects: a mixed effects model approach.
\textit{Contemp Clin Trials Commun}. 2021;22:100711.

Varadhan R, Segal JB, Boyd CM, Wu AW, Weiss CO. A framework for
the analysis of heterogeneity of treatment effect in
patient-centered outcomes research. \textit{J Clin Epidemiol}.
2013;66(8):818-825.

Wang Y, Schork NJ. Power and design issues in crossover-based
N-of-1 clinical trials with fixed data collection periods.
\textit{Healthcare (Basel)}. 2019;7(3):84.

Zucker DR, Schmid CH, McIntosh MW, D'Agostino RB, Selker HP,
Lau J. Combining single patient (N-of-1) trials to estimate
population treatment effects and to evaluate individual patient
responses to treatment. \textit{J Clin Epidemiol}. 1997;50(4):401-410.

\vfill
\noindent\rule{\textwidth}{0.4pt}
{\footnotesize
Rendered on 2026-04-15 at 10:48 PDT.\\
Source: \textasciitilde/prj/alz/10-pmsimstats-ng/pmsimstats-ng/docs/02-dgp-mean-moderation-vs-mvn.tex}
\end{document}
